\documentclass{article}
\usepackage[utf8]{inputenc}
\usepackage[english, swedish]{babel}
\usepackage{amsmath}
\usepackage{listings}
\usepackage{graphicx}
\usepackage{array}
\usepackage{longtable}
\usepackage{multirow}
\usepackage[a4paper,pdftex,bottom=20mm, width=160mm]{geometry}
\setlength{\parindent}{0pt}
\setlength{\parskip}{1em}
\usepackage{ragged2e}
\newcolumntype{L}[1]{>{\RaggedRight\arraybackslash}p{#1}}

\begin{document}

\title{Testspecifikation: Journallista}
\author{Kandidatprojekt VT 2025}
\date{\today}

\maketitle

\section{Introduktion}
Detta dokument specificerar testfall för journallistan (Lista - Översiktlig K1.1-1) i IVPJ.

\section{Testmiljö}
\subsection{Verktyg}
\begin{itemize}
    \item Pytest med Playwright
    \item Mock-data för journaler
\end{itemize}

\section{Implementerade Testfall}

\subsection{Rendering och Layout}
\begin{longtable}{|L{1.5cm}|L{3.2cm}|L{4.5cm}|L{4.5cm}|L{2.2cm}|}
\hline
\textbf{Nr} & \textbf{Testfall} & \textbf{Indata} & \textbf{Förväntat resultat} & \textbf{Krav} \\
\hline
1.1 & Lista existerar & Navigera till basvyn & Journallistan renderas och är synlig & K1.1-1 \\
\hline
1.2 & Kronologisk ordning & Lista med journaler från olika datum & Journaler visas sorterade med nyaste först & K1.2-5 \\
\hline
1.3 & Metadata-visning & Journal med titel, författare och enhet & All metadata visas korrekt för varje journal & K1.1-1 \\
\hline
1.4 & Tom lista & Ingen journaldata & Kontrollerar att inga enhetsgrupper visas när listan är tom & K1.1-1 \\
\hline
1.6 & Gruppering & Lista med journaler från samma vårdenhet & Journaler grupperas efter vårdenhet & K1.2-4 \\
\hline
1.7 & Kollapsa grupp & Klick på gruppens rubrik & Gruppen kollapsar och visar endast rubrik & K1.2-4 \\
\hline
1.8 & Expandera grupp & Klick på kollapsad grupprubrik & Gruppen expanderar och visar alla journaler & K1.2-4 \\
\hline
\end{longtable}

\subsection{Interaktion}
\begin{longtable}{|L{1.5cm}|L{3.2cm}|L{4.5cm}|L{4.5cm}|L{2.2cm}|}
\hline
\textbf{Nr} & \textbf{Testfall} & \textbf{Indata} & \textbf{Förväntat resultat} & \textbf{Krav} \\
\hline
2.1 & Markera journal & Klick på en journal & Journalen markeras med blå kant och bakgrund & K1.3-1 \\
\hline
2.2 & Avmarkera journal & Klick utanför markerad journal & Markeringen försvinner & K1.3-1 \\
\hline
2.3 & Markera flera journaler & Ctrl+klick på flera journaler & Flera journaler kan markeras samtidigt & K1.3-2 \\
\hline
\end{longtable}

\end{document} 